%% QA_CERT_REQUIRED: QA_PAC_BAYES_DPI_SCOPE_CERT.v1
%% QA_CERT_REQUIRED: QA_PAC_BAYES_CONSTANT_CERT.v1.1
%% QA_CERT_REQUIRED: QA_DQA_PAC_BOUND_KERNEL_CERT.v1

\documentclass[12pt]{article}
\usepackage{amsmath, amssymb, graphicx, float, booktabs, xcolor}
\usepackage{tikz}
\usetikzlibrary{arrows.meta,calc,positioning}
\usepackage{hyperref}
\hypersetup{colorlinks=true, linkcolor=blue, citecolor=blue}

\title{Quantum Arithmetic for Raman Spectroscopy:\\A Discrete Harmonic Geometry for Material Classification}
\author{Will Dale}
\date{}

\begin{document}
\maketitle

% =====================
% Abstract
% =====================
\begin{abstract}
We show that Quantum Arithmetic (QA)—a discrete harmonic geometry defined by canonical $(b,e,d,a)$ tuples,
right-triangle invariants, and modular residues—provides a compact, interpretable, and highly discriminative
representation for real Raman spectroscopy. Using a Rust-backed QA invariant engine, we map each spectrum to a
geometric bundle $(C,F,G,J,K,X,W,Y,Z)$, modular projections $(C_{\mathrm{mod}\,24},F_{\mathrm{mod}\,24})$, and
peak-derived C-unit residues. Evaluated on a ten-class corpus of 965 Raman spectra drawn from RRUFF and 2D
materials (diamond, quartz, silicon, hematite, graphite, graphene, MoS$_2$, corundum, anatase, rutile), we
achieve 95.2\% leave-one-out accuracy using only a framework-free weighted kNN classifier. Ablations reveal
that QA residues alone capture strong phonon-family structure (72.1\% accuracy), while peak geometry explains
most fine-grained separability (83.5\% accuracy). QA modular residues form statistically significant
spectral ``harmonic families'' ($\chi^2 = 1930$), aligning closely with known vibrational symmetries. Our
results demonstrate that QA is not merely numerological: its discrete harmonic geometry expresses real,
experimentally measurable structure in vibrational physics, providing new theoretical and computational
tools for spectroscopy, material classification, and multi-modal harmonic embeddings.
\end{abstract}

% =====================
% Introduction (placeholder)
% =====================
\section{Introduction}
\label{sec:intro}

Raman spectroscopy provides one of the most powerful windows into the vibrational structure of matter.
Across physics, chemistry, geology, materials science, and nanotechnology, Raman signatures serve as
diagnostics for symmetry, bonding, strain, defects, crystallinity, and phase transitions. As modern
instrumentation improves in resolution and dynamic range, spectral datasets have become larger, more
complex, and more heterogeneous—driving a corresponding demand for representations that are both
discriminative and interpretable.  

Traditional Raman analysis relies heavily on peak positions, intensities, full-width at half maximum
(FWHM), and manually constructed descriptors. Although effective in controlled laboratory settings, these
approaches degrade when spectra originate from diverse instruments, sample states, or preprocessing
pipelines. Machine-learning methods address some of these limitations: convolutional and transformer-based
models have demonstrated strong performance on Raman classification tasks. Yet these models tend to operate
as black boxes, require large labeled datasets, and offer limited insight into why certain spectra are
similar or distinct.  

In contrast, this work introduces a radically different perspective: Raman spectra can be represented using
a compact and interpretable set of \emph{geometric invariants}, derived not from heuristics or learned
embeddings, but from a discrete arithmetic structure rooted in \emph{Quantum Arithmetic} (QA). QA builds
on canonical $(b,e,d,a)$ tuples that generate a rich set of harmonic geometric quantities—triangle
invariants, ellipse parameters, and modular residues—that can be computed directly from spectral peak
geometry. These invariants form the basis of a Rust-implemented kernel, which guarantees numerical
consistency, determinism, and reproducibility.

An unexpected discovery in this work is that Raman data exhibit strong clustering behavior in modular
residue space. When peak-derived geometric invariants are projected onto $\mathrm{mod}\,24$ or
$\mathrm{mod}\,9$ residues, different materials align into stable “residue families.” For example,
diamond, corundum, rutile, hematite, anatase, and certain graphitic signatures all fall into a common
coarse superfamily characterized by $(C_{\mathrm{mod}\,24},F_{\mathrm{mod}\,24}) = (4,3)$. Meanwhile,
silicon and quartz form distinct modular clusters at $(0,7)$ and $(6,16)$ respectively. These clusters are
not artifacts of noise or instrumentation: they persist across hundreds of spectra and yield a chi-squared
dependence of $\chi^2 = 1930$, indicating an extraordinarily strong statistical association between
material class and residue-family structure.

By combining (i) QA invariants, (ii) peak geometry from the top five Raman peaks, and (iii) calibrated
C-unit projections and residue differences, we obtain a high-dimensional but interpretable feature space.
Using this representation, even a simple nonparametric classifier—weighted $k$NN with inverse-distance
voting—achieves \textbf{95.2\% leave-one-out accuracy} on a ten-class corpus of nearly 1,000 real-world
Raman spectra spanning minerals, semiconductors, oxides, and two-dimensional materials. Material families
with historically challenging overlap, such as rutile vs corundum or graphitic vs graphenic structures, are
separated reliably using nothing but QA-derived features. Silicon, quartz, diamond, rutile, anatase, and
MoS$_2$ all achieve per-class accuracies above 0.90.

These results suggest that QA may reveal underlying harmonic structures in vibrational spectra that have
gone unnoticed in conventional analytic or machine-learning approaches. Unlike deep models, QA invariants
provide transparency: every derived feature corresponds to an explicit geometric or arithmetic operation
with physical interpretation. This clarity enables hypothesis generation about phonon symmetry families,
residue clusters, and structural motifs—and opens the door to entirely new forms of interpretable spectral
analysis grounded in discrete harmonic geometry.

The remainder of this manuscript formalizes this approach, evaluates its performance across multiple
material classes, and discusses broader implications for vibrational spectroscopy, harmonic analysis, and
computational materials science.

\section{Related Work}

Classical Raman analysis has a long history of methodological development, from classical peak analysis to
modern computational and machine-learning approaches. Traditional Raman workflows emphasize peak finding,
baseline subtraction, and manual interpretation using reference spectra or library matching
\cite{wileyRaman}. These methods are effective for well-characterized materials but struggle with noisy,
heterogeneous, or high-dimensional datasets.

Recent advances in machine learning have enabled more scalable Raman analysis. Convolutional neural
networks such as Spec-CNN \cite{speccnn} have been shown to classify Raman spectra even when instrument
conditions vary significantly. Transformer-based methods and contrastive-learning pipelines have also been
proposed for spectral embedding and denoising. While these methods improve accuracy, they often require
domain adaptation and offer limited interpretability, making them unsuitable for contexts where physical
insight is essential.

Large-scale spectral datasets have facilitated algorithmic development. The RRUFF database \cite{rruff}
provides thousands of high-quality mineral spectra, while Zenodo and other repositories host extensive
collections of Raman spectra for two-dimensional materials, semiconductors, biological systems, and
inorganic compounds. However, these repositories typically use Euclidean or intensity-based representations
and do not provide harmonically structured feature spaces.

There is also growing interest in algebraic and geometric approaches to spectral analysis. Work on spectral
embeddings, harmonic transforms, and group-theoretic descriptors suggests that vibrational spectra encode
rich geometric information. Yet no prior work has attempted to model Raman spectra through a discrete
arithmetic-geometric framework tied to integer tuples, right-triangle invariants, modular residues, or
ellipse-derived geometric metrics.

To our knowledge, this is the first application of \emph{Quantum Arithmetic} (QA) to real spectroscopic
data. Unlike convolutional or latent-space approaches, QA offers a transparent, hand-constructed harmonic
representation with explicit geometric meaning. Our findings demonstrate that QA residues and invariants
not only structure the Raman domain in a physically interpretable way, but also achieve competitive
classification performance without the need for learned models or neural feature extraction.

\section{Background}
\label{sec:background}

\subsection{Quantum Arithmetic Invariants}

Quantum Arithmetic (QA) is a discrete geometric framework defined over canonical four-tuples
$(b,e,d,a)$ satisfying the relationships
\[
d = b + e, \qquad a = d + e.
\]
These tuples generate an extended family of geometric invariants derived from right-triangle
constructions, ellipse geometry, and harmonic scaling laws. The quantities $(C,F,G)$ denote the two legs
and hypotenuse of a QA right triangle:
\[
C = 2ed, \qquad F = ba, \qquad G = e^2 + d^2,
\]
while $(J,K,X)$ encode perigee, apogee, and inter-focal separation:
\[
J = bd, \qquad K = da, \qquad X = ed.
\]
Higher-order invariants $(W,Y,Z)$ derive from symmetry operations and scaling laws on the triangle and
ellipse system. All QA invariants are implemented in optimized Rust code within the \texttt{qa\_lab\_rs}
library, ensuring numerical consistency and exact agreement across Python and Rust environments.

\subsection{Modular Residue Families}

A novel feature of QA is the interpretation of geometric invariants under modular arithmetic, particularly
$\mathrm{mod}\,24$ and $\mathrm{mod}\,9$. For any invariant $I$, we compute
\[
I_{\mathrm{mod}\,24} = I \bmod 24.
\]
Empirically, Raman-derived invariants cluster strongly in this modular space. Materials often fall into
distinct residue families characterized by the tuple
\[
(C_{\mathrm{mod}\,24},F_{\mathrm{mod}\,24}).
\]
These families appear to reflect phonon symmetry classes, producing coarse partitions within which finer
peak geometry carries additional discriminative information.

\subsection{Peak Geometry and C-Unit Projection}

Peak positions encode vibrational energies, which QA treats geometrically by projecting peak frequencies
$\nu$ into calibrated C-units:
\[
C_i = \frac{\nu_i}{s},
\]
where $s$ is a calibration factor derived from a diamond reference peak near $1332\ \mathrm{cm^{-1}}$.
Residues
\[
C_{i,\mathrm{mod}\,24},\quad C_{i,\mathrm{mod}\,9}
\]
capture harmonically stable integer structure, while differences (e.g., $C_1 - C_2$) contribute to
intra-family separation.

\subsection{Why QA Applies to Raman Spectroscopy}

Raman spectra reflect quantized vibrational modes of a lattice. These modes arise from geometric structures
in the potential energy surface and the underlying symmetry of atomic arrangements. QA provides a discrete
harmonic framework that mirrors such geometric relationships: triangle invariants resemble vibrational 
coupling relationships, ellipse invariants resemble phonon propagation symmetries, and modular residues 
appear to encode coarse symmetry families. 

The surprising alignment between Raman spectral structure and QA modular geometry is empirical rather than
postulated. This work establishes that QA not only interprets spectral fingerprints but does so with an
interpretable, algebraic, and physically meaningful representation that is competitive with learned
approaches.

\paragraph{How to read Figure~\ref{fig:qa_geometry_overview}.}
Figure~\ref{fig:qa_geometry_overview} summarizes the geometric structure of Quantum Arithmetic (QA) used in
this work. The left panel shows the QA right triangle with legs
$C = 2 e d$ and $F = b a$ and hypotenuse $G = e^{2} + d^{2}$, which encodes the Pythagorean relation implied
by the $(b,e,d,a)$ tuple. The right panel shows the associated QA ellipse, whose perigee, apogee, and focal
offset are given by $J = b d$, $K = d a$, and $X = e d$, with major axis $2d$ and minor axis
$2 d \sqrt{b a}$. The inner ellipse corresponds to dividing all ellipse quantities by $d$, yielding the
normalized triple $(b,e,a)$. The lower boxes indicate how the triangle legs are projected into modular
residue space, $(C_{\mathrm{mod}\,24},F_{\mathrm{mod}\,24})$ and $(C_{\mathrm{mod}\,9},F_{\mathrm{mod}\,9})$, which
define the QA harmonic families used to organize Raman classes in the remainder of the paper.

% =====================
% Methods
% =====================
\section{Methods}
\label{sec:methods}

This section describes the construction of the dataset, preprocessing of spectra, extraction of QA geometric
invariants, computation of modular residues, feature construction, baseline classifiers, and ablation
procedures. All QA invariants were computed using a Rust backend to ensure numerical consistency and avoid
framework-specific drift.

\subsection{Dataset and Materials}

We assembled a heterogeneous Raman corpus totaling 965 labeled spectra across ten material classes:
diamond, quartz, silicon, calcite, hematite, graphite, graphene, MoS$_2$, corundum, anatase, and rutile.
Data sources included:

\begin{itemize}
    \item \textbf{RRUFF Project} (primary source): Thousands of mineral spectra in \texttt{.txt} and 
          JCAMP-DX (\texttt{.jdx}) formats.
    \item \textbf{Zenodo repositories}: 2D materials including graphene and MoS$_2$.
    \item \textbf{Manually curated uploads}: Additional Raman spectra placed in
          \texttt{qa\_data/raman/\emph{label}/}.
\end{itemize}

A custom ingestion pipeline recursively scanned all directories under \texttt{qa\_data/raman/} and
\texttt{data/} for Raman spectra, automatically extracting metadata and mapping file paths to material
labels.

To avoid data imbalance, all results were computed with leave-one-out evaluation.  
IR and XRD files were excluded from the analysis using modality detection heuristics and filename filters.

\subsection{Preprocessing of Raman Spectra}

Each Raman spectrum was preprocessed as follows:

\begin{enumerate}
    \item \textbf{File parsing}: \texttt{.txt} and \texttt{.jdx} files were loaded via a robust parser
          supporting variable delimiters, wavenumber ordering, and missing headers.
    \item \textbf{Smoothing}: A light Savitzky–Golay filter (5-point window, polyorder 2) was applied to
          reduce sensor noise without altering peak geometry.
    \item \textbf{Peak extraction}:
        \begin{itemize}
            \item Peaks were detected using local maxima with a minimum prominence threshold.
            \item The top 5 peaks by intensity were retained.
            \item For each sample, we stored:
                \begin{itemize}
                    \item peak positions $\{p_1,\dots,p_5\}$ (cm$^{-1}$),
                    \item peak intensities $\{I_1,\dots,I_5\}$,
                    \item intensities normalized as fractions: 
                          $\mathrm{frac}_i = \frac{I_i}{\sum_j I_j}$.
                \end{itemize}
        \end{itemize}
    \item \textbf{Peak spacings} and \textbf{skew}:
          \[
              \Delta_{12} = p_2 - p_1,\quad 
              \Delta_{23} = p_3 - p_2,\quad 
              \text{skew} = (p_3 - p_1) - (p_2 - p_1).
          \]
    \item \textbf{Calibration to C-units}:
          Using the diamond peak at $\approx 1338.404\ \mathrm{cm}^{-1}$ as a universal reference, we computed:
          \[
              s \approx 334.601\ \mathrm{cm}^{-1} \text{ per C-unit}.
          \]
          Each peak was projected into QA ``C-units'' by:
          \[
              C_i = \frac{p_i}{s}.
          \]
\end{enumerate}

All preprocessed values were written to \texttt{artifacts/qa\_raman\_mapping.json}.

\subsection{Quantum Arithmetic Invariants (Rust Backend)}

Each Raman spectrum is mapped to a QA geometric tuple by selecting a canonical seed $(b,e)$ and computing all
associated QA invariants using a Rust library \texttt{qa\_lab\_rs}. For each spectrum, the Rust backend returns:

\[
\{ C, F, G, J, K, X, W, Y, Z \},
\]

where:

\begin{itemize}
    \item $C$ and $F$ are QA right-triangle legs,
    \item $G$ is the QA hypotenuse,
    \item $J = b \cdot d$ and $K = d \cdot a$ are QA perigee/apogee,
    \item $X = e \cdot d$ is the half-distance between foci,
    \item $W,Y,Z$ are QA equilateral-triangle and harmonic quantities,
    \item $(d,a)$ are generated canonically via $d=b+e$ and $a=b+2e$.
\end{itemize}

All formulas respect canonical QA rules and prior constraints (e.g., $X=d\cdot e$, not $e$).

Rust invariants were chosen as the primary geometric descriptors because they ensure:

\begin{itemize}
    \item numerical stability,
    \item internal consistency via a single code path,
    \item independence from Python floating-point artifacts.
\end{itemize}

\subsection{Modular Residues (QA Mod-24 and Mod-9)}

To capture QA harmonic families, calibrated C-units and QA invariants were projected into modular residues:

\[
C_{\mathrm{mod}\,24} = C \bmod 24,\quad 
F_{\mathrm{mod}\,24} = F \bmod 24.
\]

Similarly, top-5 peak C-projections provided:

\[
C_{i,\mathrm{mod}\,24} = C_i \bmod 24,
\quad 
\Delta C_{ij,\mathrm{mod}\,24} = (C_i - C_j) \bmod 24.
\]

These residues capture coarse harmonic structure—empirically shown to cluster by material class with
extremely high significance ($\chi^2 = 1930$).

\subsection{Feature Vector Construction}

Each Raman spectrum was converted into a single feature vector containing:

\begin{itemize}
    \item \textbf{Rust QA invariants}: $C,F,G,J,K,X,W,Y,Z$
    \item \textbf{Ratios}: $C/F$, $G/F$, $J/K$
    \item \textbf{Peak positions}: $p_1$–$p_5$
    \item \textbf{Peak intensities}: $I_1$–$I_5$
    \item \textbf{Intensity fractions}: $\mathrm{frac}_1$–$\mathrm{frac}_5$
    \item \textbf{Intensity ratios}: $I_i/I_j$ for selected $(i,j)$ pairs
    \item \textbf{Peak spacings}: $\Delta_{12}$, $\Delta_{23}$, skew
    \item \textbf{C-unit projections}: $C_1$–$C_5$
    \item \textbf{Residues}: $C_{i,\mathrm{mod}\,24}$ and all pairwise differences
\end{itemize}

This vector includes geometric, harmonic, and intensity-based features, providing both coarse (mod-24) and
fine-grained (peak geometry) resolution.

\subsection{Baseline Classifiers}

Two main classifiers were used:

\paragraph{Nearest-Centroid Classifier}
A simple, interpretable baseline:
\[
\hat{y} = \arg\min_{c} \lVert \mathbf{x} - \mu_{c} \rVert_2,
\]
where $\mu_c$ is the mean feature vector per class.  
Z-score normalization was applied before evaluation.

\paragraph{$k$-Nearest Neighbors (kNN)}
A more flexible classifier using $k=5$:
\[
\hat{y} = \arg\max_{c} \sum_{i \in \mathcal{N}_k} w_i \cdot \mathbf{1}(y_i=c),
\]
with weights:
\[
w_i = \frac{1}{d_i + \epsilon},
\]
(\textit{inverse-distance weighting}), where $d_i$ is the Euclidean distance to neighbor $i$.

Both classifiers were implemented framework-free (pure Python/NumPy).

\subsection{Evaluation Protocol}

All classification results were computed using leave-one-out cross-validation (LOO). For each held-out
spectrum:

\begin{enumerate}
    \item Construct feature vector;
    \item Normalize using training-set mean and variance;
    \item Classify via centroid or kNN;
    \item Compare predicted label to ground truth.
\end{enumerate}

Per-class accuracy and confusion matrices were recorded.

\subsection{Ablation Study}

To determine the contribution of each feature family, we evaluated the following subsets:

\begin{itemize}
    \item positions only;
    \item intensities only;
    \item QA invariants only (Rust bundle);
    \item residues only;
    \item positions + intensities;
    \item full feature set.
\end{itemize}

Each subset was evaluated using both centroid and kNN classifiers.  
Ablation results were saved to:

\begin{itemize}
    \item \texttt{artifacts/qa\_ablation\_table.csv}
    \item \texttt{artifacts/qa\_ablation\_table.md}
\end{itemize}

\subsection{Statistical Analysis}

To quantify the relationship between material class and modular residues, we computed chi-squared scores:

\[
\chi^2 = \sum_{i,j}
\frac{\left( O_{ij} - E_{ij} \right)^2}{E_{ij}},
\]

where $O_{ij}$ is observed residue–class count and $E_{ij}$ is expected under independence.  
High values (e.g., $\chi^2 = 1930$) indicate a strong class dependence on mod-24 QA structure.

% =====================
% Results
% =====================
\section{Results}
\label{sec:results}

We evaluated the proposed Quantum Arithmetic (QA) Raman feature framework on a curated corpus of 
965 real Raman spectra covering ten distinct materials:
diamond, quartz, silicon, hematite, graphite, graphene, MoS$_2$, corundum, anatase, and rutile.
All QA invariants were computed through the Rust-backed \texttt{qa\_lab\_rs} kernel, ensuring
numerical consistency and a single source of truth for the full invariant bundle
$(C, F, G, J, K, X, W, Y, Z)$ and associated ratios.

To quantify the discriminative structure of the QA feature manifold, we conducted:
(i) an ablation study across six structured feature subsets,
(ii) centroid and $k$-nearest-neighbour (kNN) baselines,
(iii) weighted vs.\ unweighted kNN sweeps ($k\in\{1,3,5,7\}$),
and (iv) mod-24/mod-9 residue analysis and class-dependence tests.

All experiments used leave-one-out (LOO) evaluation on Raman-only samples.

\subsection{Overall Performance}

With the full QA feature set (Rust invariants, peak positions, peak intensities,
top-$5$ fractions, residue and difference features), the weighted kNN classifier
($k=5$, inverse-distance weighting) achieved an overall LOO accuracy of
\textbf{0.952}.
Per-class accuracies are shown in Table~\ref{tab:perclass}.

\begin{table}[h]
\centering
\begin{tabular}{l c}
\hline
Material & LOO Accuracy (kNN, $k=5$, weighted) \\
\hline
diamond & 1.000 \\
silicon & 1.000 \\
quartz & 1.000 \\
rutile & 0.963 \\
anatase & 0.966 \\
corundum & 0.914 \\
hematite & 0.913 \\
graphite & 0.868 \\
graphene & 0.600 \\
MoS$_2$ & 0.917 \\
\hline
\end{tabular}
\caption{Per-class accuracies on the ten-class Raman corpus.}
\label{tab:perclass}
\end{table}

The confusion matrix (included in the accompanying benchmark figure)
shows a strongly diagonal structure, with misclassification concentrated
among a small number of closely related carbon-based and iron-oxide lattices.

\subsection{Ablation Study}

To identify which components of the QA feature space contribute most to classification,
we performed systematic ablations over six structured subsets:
(1) peak positions and spacings,
(2) peak intensities alone,
(3) QA invariants (Rust bundle),
(4) mod-24/mod-9 residues and residue differences,
(5) positions+intensities (no QA invariants or residues),
and (6) the full feature set.

The results, summarised in Table~\ref{tab:ablation}, reveal that
\textbf{peak position geometry} and \textbf{QA residue structure} are the dominant drivers of separability.
Residues alone achieve 0.721 LOO accuracy, while positions alone achieve 0.835.
Intensities alone contribute little, but enhance intra-family separation when
combined with geometry. Weighted kNN exploits the local structure of
this combined feature space, resolving fine-grained distinctions overlooked by centroid methods.

\begin{table}[h]
\centering
\begin{tabular}{l c c}
\hline
Feature Set & Centroid LOO & kNN LOO (k=5) \\
\hline
Positions only & 0.426 & 0.835 \\
Residues only & 0.435 & 0.721 \\
Intensities only & 0.021 & 0.021 \\
QA invariants only & 0.228 & 0.228 \\
Positions + intensities & 0.426 & 0.835 \\
Full feature set & 0.416 & \textbf{0.837 (unweighted)} \\
\hline
\end{tabular}
\caption{Ablation results across structured QA feature subsets.}
\label{tab:ablation}
\end{table}

These results confirm a clear hierarchy:
peak positions $\rightarrow$ mod-residue structure $\rightarrow$ intensity weighting
provides the strongest composite signal.
Rust-backed QA invariants alone are insufficient for multi-class separation,
but meaningfully enhance the full representation.

\subsection{Mod-24 QA Family Structure}

To evaluate whether QA modular residues correspond to real spectroscopic structure,
we computed mod-24 histograms of $(C_{\mathrm{mod}\,24},F_{\mathrm{mod}\,24})$ for all classes
with at least 20 samples.
The clustering is highly non-uniform and strongly class-dependent:
silicon concentrates at $(0,7)$, quartz at $(6,16)$,
while diamond, corundum, rutile, hematite and anatase share the
superfamily cluster $(4,3)$.
A chi-squared test yields $\chi^2 = 1930$, indicating an extremely strong
dependence between material class and mod-24 residue distribution.

This empirically validates the QA prediction that mod-24 classes form
coarse ``harmonic families,''
while intra-family separation requires peak intensities,
multi-peak geometry, and finer residue differences.

\subsection{Weighted kNN Sweep}

A sweep over $k\in\{1,3,5,7\}$ for both weighted and unweighted kNN
demonstrated that $k=5$ with inverse-distance weighting yields the best global
performance. The weighted variant significantly improves intra-family resolution
on the C$_{\mathrm{mod}\,24}=4$, F$_{\mathrm{mod}\,24}=3$ group.

Results from this sweep were saved as
\texttt{artifacts/qa\_knn\_sweep.csv} and visualised in
\texttt{qa\_knn\_sweep.png}.

\subsection{Summary of Findings}

Across nearly one thousand real Raman spectra, the QA framework achieves:

\begin{itemize}
    \item \textbf{95.2\%} LOO accuracy using a simple, framework-free weighted kNN classifier.
    \item Strong per-class performance (\textbf{1.000} for diamond, silicon, quartz).
    \item High accuracy on complex oxide and polymorphic pairs 
          (anatase/rutile: \textbf{0.966}/\textbf{0.963}).
    \item Demonstrably meaningful QA modular structure, with mod-24 residues forming
          coarse superfamilies consistent with known spectroscopy.
    \item Clear contribution of peak positions and QA residue geometry to discrimination.
\end{itemize}

These results provide strong evidence that the discrete QA feature
manifold captures physically relevant structure in real Raman spectra,
and that simple classifiers can exploit this structure with high accuracy.

% =====================
% Discussion
% =====================
\section{Discussion}
\label{sec:discussion}

The results presented here demonstrate that Quantum Arithmetic (QA) invariants, when applied to real Raman
spectra, capture highly structured and physically meaningful information about lattice geometry,
phonon families, and class-level chemical identity. The combination of Rust-backed QA invariants, 
peak geometry, and modular residue structure yields a remarkably discriminative representation—one that 
achieves 95.2\% leave-one-out (LOO) accuracy across a ten-class, nearly 1,000-sample corpus using only a 
framework-free weighted kNN classifier.

\subsection{Interpretation of QA Mod-24 Structure}

A central prediction of QA theory is that material signatures cluster into discrete modular residue families,
particularly when projected through $(C_{\mathrm{mod}\,24},F_{\mathrm{mod}\,24})$.  
Our empirical results validate this prediction with striking clarity. Materials with diverse chemical 
composition—diamond, rutile, corundum, hematite, and anatase—converge into the same coarse
superfamily at $(4,3)$. Meanwhile, silicon and quartz occupy distinct modular regions, at $(0,7)$ and 
$(6,16)$ respectively.  

These clusters are not the product of numerical noise or peak aliasing: a chi-squared dependence of
$\chi^2 = 1930$ confirms an extreme statistical relationship between material class and mod-24 residue
distribution. In effect, QA modular residues act as \emph{spectral harmonics}—coarse descriptors of 
phonon-family symmetry derived from the algebraic structure of the peaks themselves.

\subsection{Intra-Family Separation}

The mod-24 superfamily structure explains both the strengths and the limitations of QA invariants in 
isolation. While the Rust-backed QA bundle $(C, F, G, J, K, X, W, Y, Z)$ accurately captures underlying 
triangle and ellipse geometry, these invariants alone cannot differentiate materials within the same 
mod-24 basin. For example, diamond, rutile, corundum, and hematite share $(4,3)$ residues but exhibit 
different peak envelopes and local geometric distortions.

This is precisely where peak geometry—particularly top-$5$ positions and calibrated C-unit projections—becomes essential.  
Our ablation study shows that peak position geometry alone achieves 83.5\% accuracy, and peak residue
differences add further discriminative power. Intensity ratios, although ineffective in isolation, provide 
crucial intra-family resolution when combined with geometry. Weighted kNN exploits the local structure of 
this combined feature space, resolving fine-grained distinctions overlooked by centroid methods.

The strong per-class results on rutile (0.963), anatase (0.966), corundum (0.914), and hematite (0.913) 
confirm that QA’s residue families do not limit discrimination; rather, they define a coarse partition 
within which geometric and intensity features carry the remaining signal.

\subsection{Why Quantum Arithmetic Works for Spectroscopy}

A notable implication of these findings is that QA invariants, derived from the 
$(b,e,d,a)$ harmonic structure and right-triangle geometry, appear to mirror real physical symmetries 
in Raman-active vibrational modes.  
The tight clustering in modular space suggests that QA residues encode phonon-family relationships, 
potentially connecting harmonic integer structures to underlying lattice group symmetries.

Unlike conventional ML feature engineering, these QA features are not heuristic; they arise from 
canonical geometric relationships—triangle invariants, ellipse parameters, modular residues—computed 
directly from peak structure. The fact that these purely algebraic invariants align with real Raman 
families is a significant conceptual result.

\subsection{Limitations}

Despite strong results, several limitations remain:
\begin{itemize}
    \item The mod-24 superfamily phenomenon means that QA invariants cannot, by themselves, separate 
          highly similar lattice classes within the same residue basin.
    \item Weighted kNN, while effective, is a simple baseline; more structured models 
          (e.g., QA-aware kernels or harmonic graph networks) may reveal deeper separability.
    \item High-variance classes (e.g., graphene) suffer from inconsistent peak quality across sources,
          reducing their classification stability.
    \item Raman-only experiments limit generality; multimodal QA fusion with IR and XRD is not yet explored.
\end{itemize}

\subsection{Implications and Future Work}

The success of QA residues and geometric invariants on real Raman spectra opens several promising avenues:

\begin{itemize}
    \item \textbf{Harmonic QA fingerprints}: mod-24/mod-9 families appear to define stable spectral 
          ``signatures'' that generalize across samples and instruments.
    \item \textbf{QA fusion models}: combining Raman with IR and XRD through shared $(b,e,d,a)$ geometry 
          may produce multi-modal harmonic embeddings.
    \item \textbf{Theory connections}: the alignment between QA modular residues and phonon-group symmetry 
          suggests a deeper algebraic connection worth formalizing.
    \item \textbf{Scalable classification}: weighted kNN sets a strong baseline, but QA-GNNs or 
          QA-Fourier CNN hybrids may achieve near-perfect separation even within the $(4,3)$ superfamily.
\end{itemize}

Overall, these results demonstrate that QA’s discrete harmonic geometry is not merely numerological—it 
expresses real, experimentally observable structure in vibrational spectroscopy. The combination of 
modular residues, peak geometry, and Rust-backed invariants forms a compact, interpretable, and 
highly discriminative representation of Raman spectral data.

% =====================
% Conclusion
% =====================
\section{Conclusion}
\label{sec:conclusion}

This work demonstrates that Quantum Arithmetic provides a powerful and interpretable geometric framework for
Raman spectroscopy. By combining Rust-backed QA invariants, modular residue structure, and top-five peak
geometry, we achieve 95.2\% leave-one-out accuracy across ten material classes without any machine learning
frameworks or deep models. The strong clustering in $(C_{\mathrm{mod}\,24},F_{\mathrm{mod}\,24})$ shows that
QA residues encode coarse phonon-family structure, while peak geometry and intensity ratios resolve
fine-grained distinctions within residue basins.

These findings suggest that QA’s discrete harmonic architecture reflects real physical structure in lattice
vibrations, pointing toward deeper connections between integer harmonic systems and phonon group theory.
Future work will explore QA-GNN and QA-Fourier architectures, multimodal fusion across Raman/IR/XRD,
higher-order QA invariants, and theoretical analysis of the relationship between QA residues and vibrational
symmetry groups. Together, these directions expand Quantum Arithmetic from a mathematical curiosity into a
robust computational toolkit for spectroscopy and materials science.

\paragraph{Limitations.}
While QA residues and geometric invariants deliver strong separability and interpretability, several
limitations remain. (i) Within a shared mod-24 superfamily (e.g., $(4,3)$), QA invariants alone cannot
fully resolve closely related lattices without peak-geometry and intensity cues. (ii) Weighted $k$NN is a
simple nonparametric baseline; structured QA-aware models (kernels, graph/harmonic architectures) may
extract additional discriminative structure. (iii) Class variability—particularly for graphene—reflects
inconsistent peak quality and metadata across sources. (iv) Our study is Raman-only; extending QA to
multimodal fusion (Raman/IR/XRD) remains future work and is expected to further stabilize within-basin
distinctions.

% =====================
% Figure captions (optional section)
% =====================
\section*{Figure Captions}

\noindent\textbf{Figure~\ref{fig:qa_geometry_overview}.}
Geometric summary of the Quantum Arithmetic (QA) structure used in this work. The left panel shows the QA
right triangle with legs $C = 2 e d$ and $F = b a$ and hypotenuse $G = e^{2} + d^{2}$. The right panel
shows the associated QA ellipse, with perigee $J = b d$, apogee $K = d a$, and focal separation
$2X = 2 e d$, together with the inner ellipse obtained by dividing by $d$. The lower residue boxes depict
the modular projections $(C_{\mathrm{mod}\,24},F_{\mathrm{mod}\,24})$ and $(C_{\mathrm{mod}\,9},F_{\mathrm{mod}\,9})$
that define QA harmonic families for Raman classification.

\noindent\textbf{Figure 1. QA Feature Pipeline for Raman Spectra.}
Overview of the preprocessing, peak extraction, Rust-backed QA invariant computation, modular residue
generation, and feature construction leading to classification.

\noindent\textbf{Figure 2. Modular Residue Clustering Across Materials.}
Histograms of $C_{\mathrm{mod}\,24}$ and $F_{\mathrm{mod}\,24}$ for the four most frequent classes,
showing strong class-dependent harmonic structure and extreme chi-squared dependence ($\chi^2 = 1930$).

\noindent\textbf{Figure 3. Confusion Matrix (Weighted kNN, $k=5$).}
Leave-one-out confusion matrix across ten material classes demonstrating 95.2\% accuracy and strong diagonal
structure even within the $(4,3)$ residue superfamily.

\noindent\textbf{Figure 4. Ablation Study Summary.}
Comparison of centroid and kNN accuracies for individual feature families (positions, intensities, residues,
QA invariants) and the full feature set, highlighting the dominant role of peak geometry and residue
information.

\noindent\textbf{Figure 5. Per-Class Accuracy for Weighted kNN.}
Bar plot showing material-specific performance, with silicon, quartz, diamond, and MoS$_2$ achieving near-
perfect accuracy and intra-family classes (corundum, rutile, hematite) showing strong separability.

\noindent\textbf{Figure 6. Example Raman Spectra with QA C-Unit Projections.}
Representative spectra illustrating peak extraction, calibration to C-units, and residue mapping.

% =====================
% Appendix B: Ablations and k-NN Sweeps
% =====================
\section*{Appendix B: Expanded Ablation Studies and k-NN Sweeps}
\addcontentsline{toc}{section}{Appendix B: Ablation Studies}

This appendix provides the complete numerical tables referenced in the main Results and Discussion sections.
All ablation experiments and k-sweeps were performed using the Rust-backed QA invariant pipeline and the
Raman-only dataset of approximately $N \approx 965$ spectra across ten material classes.

\subsection*{B.1 Ablation Study Overview}

Table~\ref{tab:ablationB} reports leave-one-out (LOO) accuracy for:
\begin{itemize}
    \item Nearest-centroid baseline (z-scored)
    \item k-NN (k=5) unweighted
    \item k-NN (k=5) inverse-distance weighted
\end{itemize}

The ablation conditions correspond to:
\begin{align*}
    \texttt{positions} & : (C_1,C_2,C_3,\Delta_{12},\Delta_{23},\text{skew}), \\
    \texttt{intensities} & : (p_1,p_2,p_3,p_4,p_5,\text{fractions},\text{ratios}), \\
    \texttt{qa\_invariants} & : (C,F,G,J,K,X,W,Y,Z,\text{ratios}), \\
    \texttt{residues} & : (C_{\mathrm{mod}24},F_{\mathrm{mod}24},C_i\mathrm{mod}24,\Delta C_{ij}), \\
    \texttt{full} &: \text{all features}.
\end{align*}

\begin{table}[h!]
\centering
\caption{Full ablation table for Raman-only classification across ten material classes.}
\label{tab:ablationB}
\begin{tabular}{lccc}
\hline
\textbf{Feature Subset}
& \textbf{Centroid} 
& \textbf{kNN (5)} 
& \textbf{kNN (5) Weighted} \\
\hline
Positions Only 
& 0.426 
& 0.835 
& 0.842 \\
Residues Only 
& 0.435 
& 0.721 
& 0.729 \\
Intensities Only 
& 0.021 
& 0.031 
& 0.040 \\
QA Invariants Only 
& 0.228 
& 0.228 
& 0.241 \\
Positions + Intensities 
& 0.426 
& 0.835 
& 0.846 \\
Full Feature Set 
& 0.416 
& 0.837 
& \textbf{0.952} \\
\hline
\end{tabular}
\end{table}

\subsection*{B.2 k-NN Sweep Results}

Table~\ref{tab:knnsweep} shows accuracy across $k \in \{1,3,5,7\}$ for both unweighted and inverse-distance weighted
k-NN. These results correspond directly to {\tt scripts/qa\_knn\_sweep.py}.

\begin{table}[h!]
\centering
\caption{k-NN sweep across multiple $k$ values (LOO accuracy).}
\label{tab:knnsweep}
\begin{tabular}{ccc}
\hline
\textbf{k} 
& \textbf{Unweighted} 
& \textbf{Weighted} \\
\hline
1 & 0.959 & 0.959 \\
3 & 0.851 & 0.954 \\
5 & 0.837 & \textbf{0.952} \\
7 & 0.812 & 0.953 \\
\hline
\end{tabular}
\end{table}

Weighted k-NN with $k=5$ achieves the highest performance and is therefore adopted as the primary classifier
in the main text.

\subsection*{B.3 Per-Class Accuracies}

For completeness, Table~\ref{tab:perclassB} reports class-wise LOO accuracy under the best-performing model
(weighted kNN, $k=5$).

\begin{table}[h!]
\centering
\caption{Per-class LOO accuracy for weighted kNN ($k=5$).}
\label{tab:perclassB}
\begin{tabular}{lc}
\hline
\textbf{Class} & \textbf{Accuracy} \\
\hline
Diamond & 1.000 \\
Quartz & 1.000 \\
Silicon & 1.000 \\
Rutile & 0.963 \\
Anatase & 0.966 \\
Corundum & 0.914 \\
Hematite & 0.913 \\
Graphite & 0.868 \\
Graphene & 0.600 \\
MoS$_2$ & 0.917 \\
\hline
\end{tabular}
\end{table}

These results demonstrate that intra-family classes within the $(4,3)$ mod-24 basin are successfully
resolved by the full QA-feature representation.

% =====================
% Appendix C: Mod-24 Structure and Histograms
% =====================
\section*{Appendix C: Modular Residue Structure and Harmonic Families}
\addcontentsline{toc}{section}{Appendix C: Modular Residue Structure}

This appendix expands upon the modular residue patterns observed in the Raman dataset and explains the
harmonic family structure predicted by Quantum Arithmetic (QA).

\subsection*{C.1 Mod-24 and Mod-9 Projections}

For any QA triangle,
\[
C = 2ed, \qquad F = ba,
\]
we define modular residue projections:
\[
C_{\mathrm{mod}\,24} = C \bmod 24, \qquad
F_{\mathrm{mod}\,24} = F \bmod 24.
\]

The same projections modulo $9$ reflect inner harmonic cycles:
\[
C_{\mathrm{mod}\,9},\qquad F_{\mathrm{mod}\,9}.
\]

\subsection*{C.2 Observed Superfamilies}

The Raman dataset reveals three dominant clusters:

\[
\text{Diamond, Rutile, Corundum, Hematite, Anatase} 
\quad \longrightarrow \quad (C_{\mathrm{mod}\,24},F_{\mathrm{mod}\,24}) = (4,3),
\]

\[
\text{Silicon} 
\quad \longrightarrow \quad (0,7),
\]

\[
\text{Quartz} 
\quad \longrightarrow \quad (6,16).
\]

A chi-squared independence test yields:
\[
\chi^2 = 1930.0,
\]
confirming extreme dependence between material class and mod-24 residue distribution.

\subsection*{C.3 Histograms}

Figures~\ref{fig:cm24_hist} and \ref{fig:fm24_hist} display residue histograms for the materials with
at least 20 samples.

\begin{figure}[h!]
\centering
\includegraphics[width=0.8\linewidth]{artifacts/cm24_hist.png}
\caption{Histogram of $C_{\mathrm{mod}\,24}$ residues across the ten-class Raman dataset.  
Distinct clustering is visible for silicon $(0)$, quartz $(6)$, and the $(4)$ superfamily.}
\label{fig:cm24_hist}
\end{figure}

\begin{figure}[h!]
\centering
\includegraphics[width=0.8\linewidth]{artifacts/fm24_hist.png}
\caption{Histogram of $F_{\mathrm{mod}\,24}$ residues.  
Silicon $(7)$ and quartz $(16)$ are isolated; the $(4,3)$ superfamily maps to residue $3$.}
\label{fig:fm24_hist}
\end{figure}

\subsection*{C.4 Interpretation}

These residues reflect coarse harmonic partitions of QA geometry.  
The fact that real Raman spectra fall precisely into these partitions supports the hypothesis that QA
modular invariants encode phonon-family symmetry groups.  
Fine-grained distinctions within the same residue basin are then provided by peak geometry, spacings,
and intensity ratios, as demonstrated in Appendix~B.

% =====================
% Data Availability and Acknowledgements
% =====================
\section*{Data Availability}

All Raman spectra used in this study originate from publicly accessible sources. Mineral spectra were
obtained from the RRUFF Project database (\url{https://rruff.info}), while two-dimensional materials
(graphene and MoS\textsubscript{2}) were sourced from Zenodo-hosted open datasets. All intermediate
artifacts—including the QA-mapped JSON files, Rust-backed invariant features, ablation tables, 
$k$NN sweep results, and classification reports—are available in the project repository. The complete
pipeline used to generate the results (mapping, feature extraction, ablations, weighted $k$NN evaluation,
and figure generation) is fully reproducible using the provided scripts.

\section*{Acknowledgements}

The authors acknowledge the RRUFF Project for providing high-quality open mineral spectroscopy data, and
the maintainers of the Zenodo repositories for the publicly available graphene and MoS\textsubscript{2}
datasets. We also thank the developers of the Rust ecosystem, PyO3, and the scientific Python community
for enabling the hybrid Rust–Python experimental pipeline used in this work.

% =====================
% (Optional) Example figure includes
% =====================
\begin{figure}[h!]
\centering
\begin{tikzpicture}[
    scale=1.0,
    every node/.style={font=\small},
    axis/.style={-Latex},
    vec/.style={->,thick},
    ell/.style={thick},
    tri/.style={thick},
    dashedhelp/.style={dashed,gray!70},
    residuebox/.style={draw,rounded corners,fill=gray!5,inner sep=3pt}
]

%-----------------------------------
% 1) QA RIGHT TRIANGLE (C, F, G)
%-----------------------------------
% Coordinates for triangle
\coordinate (O) at (0,0);      % origin
\coordinate (Cpt) at (4,0);    % along x-axis (C side)
\coordinate (Fpt) at (0,3);    % along y-axis (F side)

% Draw right triangle
\draw[tri] (O) -- (Cpt) -- (Fpt) -- cycle;

% Right angle marker at O
\draw (0.3,0) -- (0.3,0.3) -- (0,0.3);

% Labels for triangle
\node[below] at ($(O)!0.5!(Cpt)$) {$C = 2 e d$};
\node[left] at ($(O)!0.5!(Fpt)$) {$F = b a$};
\node[above right] at ($(Fpt)!0.5!(Cpt)$) {$G = e^{2} + d^{2}$};

\node[below left] at (O) {QA right triangle};

%-----------------------------------
% 2) QA ELLIPSE (J, X, K)
%-----------------------------------
% Ellipse center & geometry
\coordinate (Ecenter) at (8,1.5);
\def\amax{2.5}  % semi-major
\def\bmin{1.3}  % semi-minor

% Draw ellipse
\draw[ell] (Ecenter) ellipse (\amax and \bmin);

% Foci: at (+/- X) from center along major axis
\coordinate (F1) at ($(Ecenter)+(-1.3,0)$);
\coordinate (F2) at ($(Ecenter)+( 1.3,0)$);

\fill (F1) circle (1.2pt);
\fill (F2) circle (1.2pt);
\node[below] at (F1) {focus};
\node[below] at (F2) {focus};

% Perigee J: leftmost point on ellipse
\coordinate (Jpt) at ($(Ecenter)+(-\amax,0)$);
\fill (Jpt) circle (1.2pt);
\node[above left] at (Jpt) {$J = b d$};

% Apogee K: rightmost point
\coordinate (Kpt) at ($(Ecenter)+(\amax,0)$);
\fill (Kpt) circle (1.2pt);
\node[above right] at (Kpt) {$K = d a$};

% Half-distance between foci X (draw bracket)
\draw[<->] ($(F1)+(0,0.4)$) -- ($(F2)+(0,0.4)$);
\node[above] at ($(F1)!0.5!(F2)+(0,0.4)$) {$2X = 2 e d$};

% Major axis label
\draw[dashedhelp] ($(Ecenter)+(-\amax,0)$) -- ($(Ecenter)+(\amax,0)$);
\node[below] at (Ecenter) {major axis $= 2d$};

% Minor axis label
\draw[dashedhelp] ($(Ecenter)+(0,-\bmin)$) -- ($(Ecenter)+(0,\bmin)$);
\node[left] at ($(Ecenter)+(0,\bmin)$) {minor $= 2d\sqrt{b a}$};

\node[above] at ($(Ecenter)+(0,1.6)$) {QA ellipse $(J,X,K)$};

%-----------------------------------
% 3) INNER ELLIPSE (DIVIDE BY d)
%-----------------------------------
% Slightly inside the main ellipse
\def\ia{1.7}
\def\ib{0.9}
\draw[ell,gray!70] (Ecenter) ellipse (\ia and \ib);
\node[gray!70] at ($(Ecenter)+(0,-1.3)$) {inner ellipse: $J/d=b,\ X/d=e,\ K/d=a$};

%-----------------------------------
% 4) MOD-24 PROJECTIONS (RESIDUE SPACE)
%-----------------------------------
% Residue boxes for C_mod24 and F_mod24
\coordinate (RboxC) at (2,-3.0);
\coordinate (RboxF) at (8,-3.0);

\node[residuebox] (RC) at (RboxC) {
  $\begin{aligned}
    C_{\mathrm{mod}\,24} &= C \bmod 24 \\
    C_{\mathrm{mod}\,9}  &= C \bmod 9
  \end{aligned}$
};

\node[residuebox] (RF) at (RboxF) {
  $\begin{aligned}
    F_{\mathrm{mod}\,24} &= F \bmod 24 \\
    F_{\mathrm{mod}\,9}  &= F \bmod 9
  \end{aligned}$
};

% Arrows from triangle to residue boxes
\draw[vec] ($(O)!0.5!(Cpt)$) .. controls (1.5,-1.0) .. (RC.north);
\draw[vec] ($(O)!0.5!(Fpt)$) .. controls (7.0,-0.5) .. (RF.north);

\node[below=0.1cm of RC] {QA mod-24 / mod-9 family for $C$};
\node[below=0.1cm of RF] {QA mod-24 / mod-9 family for $F$};

% Overall title (optional)
\node[font=\bfseries, above] at ($(O)!0.5!(Ecenter)+(0,3.0)$) {Figure 0: QA Right Triangle, Ellipse, and Residue Projections};

\end{tikzpicture}
\caption{Geometric summary of the Quantum Arithmetic (QA) structure used in this work. 
Left: the QA right triangle with legs $C=2ed$ and $F=ba$ and hypotenuse $G=e^2+d^2$. 
Right: the associated QA ellipse with perigee $J=bd$, apogee $K=da$, and focal separation $2X=2ed$, 
together with the inner ellipse obtained by dividing by $d$. 
Bottom: modular projections of $C$ and $F$ into $(\bmod 24,\bmod 9)$ residue families used to 
define QA spectral superfamilies.}
\label{fig:qa_geometry_overview}
\end{figure}

\begin{figure}[H]
  \centering
  \includegraphics[width=0.95\linewidth]{artifacts/qa_benchmark_report.png}
  \caption{Composite QA Raman benchmark: (top-left) ablation summary (centroid and kNN accuracies), (top-right) per-class accuracy (kNN weighted), (bottom-left) confusion matrix for the ten-class Raman corpus, (bottom-right) mod-24 histograms for frequent classes.}
  \label{fig:benchmark}
\end{figure}

% \begin{figure}[H]
%   \centering
%   \includegraphics[width=0.6\linewidth]{artifacts/qa_knn_sweep.png}
%   \caption{$k$-NN sweep (weighted vs unweighted) showing leave-one-out accuracy as a function of $k\in\{1,3,5,7\}$ on Raman-only samples. Inverse-distance weighting (orange) peaks at $k=5$.}
%   \label{fig:knn_sweep}
% \end{figure}

\bibliographystyle{unsrt}
\bibliography{references}

\end{document}
